\documentclass[a4paper,12pt]{article}

\usepackage[T2A]{fontenc}
\usepackage[utf8]{inputenc}
\usepackage[russian]{babel}
\usepackage{amsmath,amssymb}
\usepackage{helvet}
\renewcommand{\familydefault}{\sfdefault}
\usepackage{icomma}
\usepackage[a4paper,left=20mm,right=20mm,top=19mm,bottom=20mm]{geometry}
\usepackage{microtype}
\usepackage{tikz}
\usepackage{tkz-euclide}
\usepackage{tabularx,booktabs}
\usepackage{enumitem}
\usepackage{hyperref}
\usepackage{parskip}
\usepackage{fancyhdr}
\usepackage{setspace}
\usepackage{xcolor}

\definecolor{accent}{HTML}{008080}
\definecolor{darkaccent}{HTML}{004D4D}
\definecolor{textgray}{HTML}{333333}
\definecolor{softgray}{HTML}{6B7280}
\definecolor{rulegray}{HTML}{D8DEDE}

\hypersetup{
  colorlinks=true,
  linkcolor=darkaccent,
  urlcolor=darkaccent,
  pdfauthor={Лёвин Артём Александрович},
  pdftitle={Чек-лист: медианы треугольника}
}

\onehalfspacing
\setlength{\parindent}{0pt}
\setlist[itemize]{leftmargin=1.4em,itemsep=0.25em,topsep=0.35em}
\setlist[enumerate]{leftmargin=1.7em,itemsep=0.45em,topsep=0.4em}
\raggedbottom

\pagestyle{fancy}
\fancyhf{}
\fancyhead[L]{\small\color{softgray}ОГЭ · Геометрия}
\fancyhead[R]{\small\color{softgray}Медианы треугольника}
\fancyfoot[C]{\small\color{softgray}\thepage}
\setlength{\headheight}{25pt}
\renewcommand{\headrulewidth}{0.4pt}
\renewcommand{\headrule}{%
  \hbox to\headwidth{%
    \color{rulegray}\leaders\hrule height \headrulewidth\hfill
  }%
}

\tkzSetUpLine[line width=0.9pt,color=accent]
\tkzSetUpPoint[color=darkaccent,fill=darkaccent,size=3]
\tkzSetUpLabel[color=black,font=\small]

\newcommand{\checksection}[1]{%
  \vspace{0.7em}%
  {\Large\bfseries\color{darkaccent}#1}\par
  \vspace{0.25em}
  {\color{rulegray}\rule{\linewidth}{0.6pt}}\par
  \vspace{0.5em}
}

\newcommand{\key}[1]{\textbf{\color{darkaccent}#1}}
\newcommand{\SABC}{S_{ABC}}

\begin{document}

% ============================================================
% ТИТУЛЬНЫЙ ЛИСТ
% ============================================================

\begin{titlepage}
\thispagestyle{empty}
\centering

\vspace*{28mm}

{\Large\color{softgray}Чек-лист\par}

\vspace{5mm}

{\Huge\bfseries\color{darkaccent}
Медианы треугольника\par}

\vspace{3mm}

{\LARGE\color{textgray}
площади, центр тяжести, удвоение медианы, подобие\par}

\vspace{18mm}

{\large ОГЭ · геометрия\par}

\vfill

{\large\bfseries
Лёвин Артём Александрович\par}

\vspace{2mm}

{\large
эксклюзивно для Софии Кальней\par}

\vspace{7mm}

{\normalsize \today\par}

\vspace{18mm}

\begin{tikzpicture}[x=1cm,y=1cm]
  \draw[color=accent,line width=1.1pt]
    (-7,0)
    .. controls (-4.3,0.8) and (-2.2,-0.55) .. (0,0.05)
    .. controls (2.6,0.75) and (4.5,-0.55) .. (7,0.1);
\end{tikzpicture}

\vspace*{3mm}

\end{titlepage}

% ============================================================
% ОСНОВНОЙ ЧЕК-ЛИСТ
% ============================================================

\checksection{1. Что нужно знать о медиане}

\key{Медиана треугольника} --- отрезок, соединяющий вершину
треугольника с серединой противоположной стороны.

Пусть в $\triangle ABC$ точки $A_1$, $B_1$, $C_1$ ---
середины сторон $BC$, $CA$, $AB$. Тогда
\[
AA_1,\qquad BB_1,\qquad CC_1
\]
--- медианы треугольника.

\begin{center}
\begin{tikzpicture}[scale=0.95]

  \tkzDefPoint(0,4){A}
  \tkzDefPoint(-3,0){B}
  \tkzDefPoint(4,0){C}

  \tkzDefMidPoint(B,C)
  \tkzGetPoint{A1}

  \tkzDefMidPoint(C,A)
  \tkzGetPoint{B1}

  \tkzDefMidPoint(A,B)
  \tkzGetPoint{C1}

  \tkzInterLL(A,A1)(B,B1)
  \tkzGetPoint{O}

  \tkzDrawPolygon(A,B,C)
  \tkzDrawSegments(A,A1 B,B1 C,C1)

  \tkzDrawPoints(A,B,C,A1,B1,C1,O)

  \tkzLabelPoints[above](A)
  \tkzLabelPoints[below left](B)
  \tkzLabelPoints[below right](C)
  \tkzLabelPoints[below](A1)
  \tkzLabelPoints[right](B1)
  \tkzLabelPoints[left](C1)
  \tkzLabelPoint[above right](O){$O$}

\end{tikzpicture}
\end{center}

\key{Главное свойство точки пересечения медиан.}

Все три медианы пересекаются в одной точке $O$ и делятся
этой точкой в отношении $2:1$, считая от вершины:
\[
AO:OA_1
=
BO:OB_1
=
CO:OC_1
=
2:1.
\]

Следовательно,
\[
AO=\frac23 AA_1,
\qquad
OA_1=\frac13 AA_1,
\]
и аналогично для двух других медиан.

\textbf{Контрольная мысль.}
Вся медиана состоит из трёх одинаковых долей.
От вершины до точки пересечения медиан приходится две доли,
от точки пересечения до середины стороны --- одна.

% ------------------------------------------------------------

\checksection{2. Медианы и площади}

\key{Одна медиана делит треугольник на две равновеликие части.}

Если $AA_1$ --- медиана, то
\[
S_{ABA_1}=S_{ACA_1}=\frac12\SABC.
\]

Причина: у треугольников $ABA_1$ и $ACA_1$ общая высота,
проведённая к прямой $BC$, а основания
\[
BA_1=A_1C.
\]

\key{Три медианы делят треугольник на шесть равновеликих
треугольников.}

Площадь каждого малого треугольника равна
\[
\boxed{\frac16\SABC}.
\]

Важно: эти шесть треугольников гарантированно
\emph{равны по площади}. Их равенство как геометрических фигур
из этого свойства не следует.

Полезное следствие:
\[
S_{AOB}=S_{BOC}=S_{COA}
=
\frac13\SABC,
\]
поскольку каждый из этих треугольников состоит из двух малых
треугольников площади $\frac16\SABC$.

% ------------------------------------------------------------

\checksection{3. Удвоение медианы: главное дополнительное построение}

Пусть $AD$ --- медиана $\triangle ABC$.
Продлим её за точку $D$ до точки $K$ так, чтобы
\[
AD=DK.
\]

Точка $D$ является серединой отрезка $BC$ и серединой отрезка
$AK$. Следовательно, диагонали четырёхугольника $ABKC$
делятся пополам.

Отсюда
\[
\boxed{ABKC\text{ --- параллелограмм}.}
\]

\begin{center}
\begin{tikzpicture}[scale=0.9]

  \tkzDefPoint(-2.5,0){A}
  \tkzDefPoint(1.5,0.5){B}
  \tkzDefPoint(-0.3,3.5){C}
  \tkzDefPoint(0.6,2){D}
  \tkzDefPoint(3.7,4){K}

  \tkzDrawPolygon(A,B,C)
  \tkzDrawSegments(A,K B,K C,K)

  \tkzDrawPoints(A,B,C,D,K)

  \tkzLabelPoints[below left](A)
  \tkzLabelPoints[below right](B)
  \tkzLabelPoints[above left](C)
  \tkzLabelPoints[above right](K)
  \tkzLabelPoint[right](D){$D$}

\end{tikzpicture}
\end{center}

Из свойств параллелограмма получаем
\[
AB=CK,
\qquad
AC=BK,
\]
\[
AB\parallel CK,
\qquad
AC\parallel BK.
\]

\key{Зачем это нужно.}

Удвоение медианы позволяет перейти к параллельным прямым,
равным отрезкам, подобным и прямоугольным треугольникам.
Это один из основных приёмов решения задач повышенной
сложности на медианы.

% ------------------------------------------------------------

\checksection{4. Задача-эталон: медианы равны $3$, $4$, $5$}

\textbf{Задача.}
Найдите площадь треугольника, медианы которого равны
$3$, $4$ и $5$.

Обозначим точку пересечения медиан через $O$.

Отрезок от вершины до точки $O$ составляет $\frac23$
соответствующей медианы. После удвоения подходящего отрезка
получается треугольник $OCF$ со сторонами
\[
OC=\frac23\cdot5=\frac{10}{3},
\]
\[
CF=\frac23\cdot3=2,
\]
\[
OF=\frac23\cdot4=\frac83.
\]

\begin{center}
\begin{tikzpicture}[scale=1.15]

  \tkzDefPoint(0,0){F}
  \tkzDefPoint(2,0){C}
  \tkzDefPoint(0,2.6667){O}

  \tkzDefMidPoint(O,F)
  \tkzGetPoint{B1}

  \tkzDrawPolygon(O,C,F)
  \tkzDrawSegment(C,B1)

  \tkzDrawPoints(O,C,F,B1)

  \tkzMarkRightAngle[size=0.28](C,F,O)

  \tkzLabelPoints[above](O)
  \tkzLabelPoints[right](C)
  \tkzLabelPoints[below left](F)
  \tkzLabelPoint[left](B1){$B_1$}

  \tkzLabelSegment[above right](O,C){$\frac{10}{3}$}
  \tkzLabelSegment[below](F,C){$2$}
  \tkzLabelSegment[left](O,F){$\frac83$}

\end{tikzpicture}
\end{center}

Полупериметр треугольника $OCF$:
\[
p=
\frac{
\frac{10}{3}+2+\frac83
}{2}
=4.
\]

По формуле Герона
\[
S_{OCF}
=
\sqrt{
4
\left(4-\frac{10}{3}\right)
(4-2)
\left(4-\frac83\right)
}.
\]

Получаем
\[
S_{OCF}
=
\sqrt{
4\cdot
\frac23\cdot
2\cdot
\frac43
}
=
\frac83.
\]

Отрезок $CB_1$ является медианой $\triangle OCF$, поэтому
\[
S_{OCB_1}
=
\frac12S_{OCF}
=
\frac43.
\]

Треугольник $OCB_1$ является одним из шести равновеликих
треугольников исходного треугольника. Следовательно,
\[
\SABC
=
6\cdot\frac43
=
\boxed{8}.
\]

% ------------------------------------------------------------

\checksection{5. Треугольник, построенный на медианах}

Пусть стороны некоторого нового треугольника равны трём
медианам исходного $\triangle ABC$.

Обозначим площадь такого треугольника через $S_m$.

После удвоения медианы возникает треугольник, каждая сторона
которого равна $\frac23$ соответствующей медианы исходного
треугольника.

Следовательно, этот треугольник подобен треугольнику,
построенному на медианах, с коэффициентом
\[
k=\frac23.
\]

Площади подобных треугольников относятся как квадраты
коэффициента подобия:
\[
S_{OCF}
=
\left(\frac23\right)^2S_m
=
\frac49S_m.
\]

При этом треугольник $OCF$ состоит из двух малых треугольников,
каждый из которых имеет площадь $\frac16\SABC$. Поэтому
\[
S_{OCF}
=
2\cdot\frac16\SABC
=
\frac13\SABC.
\]

Получаем
\[
\frac49S_m
=
\frac13\SABC.
\]

Отсюда
\[
\boxed{S_m=\frac34\SABC}.
\]

Обратная формула:
\[
\boxed{\SABC=\frac43S_m}.
\]

\textbf{Практический смысл.}

Если известны три медианы исходного треугольника, можно:

\begin{enumerate}
  \item построить мысленно обычный треугольник с такими сторонами;
  \item найти его площадь;
  \item умножить полученную площадь на $\frac43$.
\end{enumerate}

% ------------------------------------------------------------

\checksection{6. Подобие и прямоугольные треугольники}

В сложной геометрической задаче полезно выполнять два
автоматических поиска.

\begin{itemize}

  \item
  \key{Прямоугольные треугольники.}
  Высота создаёт прямой угол. После этого можно применять
  теорему Пифагора и другие свойства прямоугольного
  треугольника.

  \item
  \key{Подобные треугольники.}
  Особенно внимательно проверяйте вертикальные углы и углы
  при параллельных прямых, появившихся после дополнительного
  построения.

\end{itemize}

\textbf{Типовая конфигурация.}

Пусть в равнобедренном $\triangle ABC$
\[
AB=BC.
\]

Медиана $AD$ и высота $CE$ пересекаются в точке $P$.
Удвоим медиану $AD$ до точки $K$.

Тогда
\[
ABKC
\]
--- параллелограмм и
\[
CK=AB.
\]

\begin{center}
\begin{tikzpicture}[scale=0.62]

  \tkzDefPoint(0,0){A}
  \tkzDefPoint(8.75,0){B}
  \tkzDefPoint(3.5,7){C}
  \tkzDefPoint(6.125,3.5){D}
  \tkzDefPoint(3.5,0){E}
  \tkzDefPoint(3.5,2){P}
  \tkzDefPoint(12.25,7){K}

  \tkzDrawPolygon(A,B,C)
  \tkzDrawSegments(A,K C,E B,K C,K)

  \tkzDrawPoints(A,B,C,D,E,P,K)

  \tkzMarkRightAngle[size=0.33](C,E,B)

  \tkzLabelPoints[below](A,E,B)
  \tkzLabelPoints[above](C,K)
  \tkzLabelPoint[above right](D){$D$}
  \tkzLabelPoint[left](P){$P$}

\end{tikzpicture}
\end{center}

Рассмотрим треугольники
\[
\triangle APE
\quad\text{и}\quad
\triangle KPC.
\]

Они подобны:

\begin{itemize}
  \item
  \[
  \angle APE=\angle KPC
  \]
  как вертикальные;

  \item
  \[
  AE\parallel KC,
  \]
  поэтому соответствующие углы равны.
\end{itemize}

Следовательно,
\[
\frac{AE}{KC}
=
\frac{PE}{PC}.
\]

Пусть
\[
AB=BC=x.
\]

Так как $ABKC$ --- параллелограмм,
\[
KC=x.
\]

Поэтому отношение
\[
\frac{PE}{PC}
\]
позволяет выразить $AE$ через $x$, а затем найти
\[
BE=AB-AE.
\]

После этого применяется теорема Пифагора в прямоугольном
$\triangle BEC$.

\textbf{Пример.}

Пусть
\[
PC=5,
\qquad
PE=2.
\]

Тогда
\[
CE=PC+PE=7.
\]

Из подобия:
\[
\frac{AE}{x}
=
\frac25.
\]

Следовательно,
\[
AE=\frac25x,
\qquad
BE=x-\frac25x=\frac35x.
\]

В прямоугольном $\triangle BEC$:
\[
BC^2=BE^2+CE^2.
\]

Получаем
\[
x^2
=
\left(\frac35x\right)^2
+
7^2.
\]

Следовательно,
\[
x^2
=
\frac9{25}x^2+49,
\]
\[
\frac{16}{25}x^2=49,
\]
\[
x^2=\frac{1225}{16},
\]
\[
x=\frac{35}{4}.
\]

Так как $CE$ --- высота к стороне $AB$,
\[
S_{ABC}
=
\frac12\cdot AB\cdot CE.
\]

Поэтому
\[
S_{ABC}
=
\frac12
\cdot
\frac{35}{4}
\cdot
7
=
\boxed{\frac{245}{8}}.
\]

% ------------------------------------------------------------

\checksection{7. Формулы, которые должны быть на автомате}

\renewcommand{\arraystretch}{1.35}

\begin{center}
\begin{tabularx}{\linewidth}{@{}>{\bfseries}p{0.34\linewidth}X@{}}

\toprule
Ситуация & Что использовать \\
\midrule

Точка пересечения медиан
&
$AO=\frac23AA_1$,
$OA_1=\frac13AA_1$
\\

Одна медиана
&
Делит площадь исходного треугольника пополам
\\

Три медианы
&
Шесть равновеликих треугольников,
каждый площади $\frac16S$
\\

Удвоение медианы
&
Получение параллелограмма через взаимное деление
диагоналей пополам
\\

Площади подобных треугольников
&
$\displaystyle \frac{S_1}{S_2}=k^2$
\\

Треугольник на медианах
&
$S_m=\frac34S$,
\quad
$S=\frac43S_m$
\\

Три известные стороны
&
Формула Герона:
$\displaystyle
S=\sqrt{p(p-a)(p-b)(p-c)}
$
\\

Высота и основание
&
$\displaystyle S=\frac12ah_a$
\\

Прямоугольный треугольник
&
Теорема Пифагора
\\

\bottomrule

\end{tabularx}
\end{center}

% ------------------------------------------------------------

\checksection{8. Типичные ошибки}

\begin{itemize}

  \item
  Путать отношение $2:1$ с отношением $1:2$.
  Две части медианы лежат со стороны вершины.

  \item
  Записывать
  \[
  AO=\frac23OA_1.
  \]
  Верные соотношения:
  \[
  AO=2\,OA_1,
  \qquad
  AO=\frac23AA_1.
  \]

  \item
  Считать шесть малых треугольников равными как фигуры.
  Гарантируется равенство их площадей.

  \item
  После удвоения медианы использовать свойства
  параллелограмма до проверки того, что его диагонали
  делятся пополам.

  \item
  Сопоставлять стороны подобных треугольников только по
  внешнему виду чертежа. Соответствие сторон определяется
  соответствием углов.

  \item
  Забывать возвести коэффициент подобия в квадрат при
  сравнении площадей.

  \item
  Применять формулу Герона без предварительного вычисления
  полупериметра
  \[
  p=\frac{a+b+c}{2}.
  \]

  \item
  На последнем шаге забывать, площадь какого именно
  треугольника была найдена: малого, дополнительного,
  треугольника на медианах или исходного.

\end{itemize}

% ------------------------------------------------------------

\checksection{9. Универсальный алгоритм решения}

\begin{enumerate}

  \item
  Отметьте середины сторон и все данные о медианах.

  \item
  Если медианы пересекаются, сразу запишите отношение
  \[
  2:1.
  \]

  \item
  Если речь идёт о площадях, проверьте возможность
  использовать доли
  \[
  \frac12,\qquad
  \frac13,\qquad
  \frac16.
  \]

  \item
  Если дана медиана и исходная конфигурация неудобна,
  рассмотрите её удвоение.

  \item
  После дополнительного построения выпишите равные и
  параллельные стороны полученного параллелограмма.

  \item
  Найдите прямоугольные треугольники.

  \item
  Найдите пары подобных треугольников.
  В первую очередь проверяйте вертикальные углы и углы
  при параллельных прямых.

  \item
  Переведите геометрические связи в отношения,
  уравнения и формулы площадей.

  \item
  Выполните вычисления.

  \item
  Проверьте, во сколько раз найденная площадь отличается
  от требуемой.

\end{enumerate}

% ============================================================
% ТРЕНИРОВОЧНЫЙ БЛОК
% ============================================================

\newpage

\thispagestyle{fancy}

\begin{center}
{\LARGE\bfseries\color{darkaccent}
Тренировочный блок · Путь А\par}

\vspace{2mm}

{\color{softgray}
10 задач по нарастающей сложности\par}
\end{center}

\vspace{4mm}

\textbf{Средний уровень}

\begin{enumerate}[label=\textbf{\arabic*.}]

  \item
  В $\triangle ABC$ медиана
  \[
  AA_1=12.
  \]
  Точка $O$ --- точка пересечения медиан.
  Найдите $AO$ и $OA_1$.

  \item
  Площадь треугольника равна $54$.
  Проведены все три медианы.
  Найдите площадь каждого из шести малых треугольников.

  \item
  Медиана $AD$ делит $\triangle ABC$ на два треугольника.
  Известно, что
  \[
  S_{ABD}=17.
  \]
  Найдите $S_{ABC}$.

\end{enumerate}

\vspace{2mm}

\textbf{Уровень выше среднего}

\begin{enumerate}[label=\textbf{\arabic*.},resume]

  \item
  Медианы треугольника равны
  \[
  6,\qquad 8,\qquad 10.
  \]
  Найдите площадь исходного треугольника.

  \item
  Медианы треугольника равны
  \[
  5,\qquad 5,\qquad 6.
  \]
  Найдите площадь исходного треугольника.

  \item
  Площадь исходного треугольника равна $120$.
  Найдите площадь треугольника, стороны которого равны
  медианам исходного треугольника.

  \item
  В $\triangle ABC$ медиана $AD$ продолжена до точки $K$
  так, что
  \[
  AD=DK.
  \]
  Известно:
  \[
  AB=7,
  \qquad
  AC=9.
  \]
  Найдите $CK$ и $BK$.

  \item
  Площадь $\triangle ABC$ равна $90$.
  Все три медианы пересекаются в точке $O$.
  Найдите площади треугольников
  \[
  AOB,\qquad BOC,\qquad COA.
  \]

  \item
  В равнобедренном $\triangle ABC$
  \[
  AB=BC.
  \]
  Медиана $AD$ и высота $CE$ пересекаются в точке $P$.
  Известно:
  \[
  PC=6,
  \qquad
  PE=2.
  \]
  Найдите площадь $\triangle ABC$.

\end{enumerate}

\vspace{2mm}

\textbf{Сложная задача}

\begin{enumerate}[label=\textbf{\arabic*.},resume]

  \item
  В равнобедренном $\triangle ABC$
  \[
  AB=BC.
  \]
  Медиана $AD$ к стороне $BC$ и высота $CE$ к стороне
  $AB$ пересекаются в точке $P$.

  Известно:
  \[
  PC=9,
  \qquad
  PE=3.
  \]

  Найдите площадь $\triangle ABC$.

\end{enumerate}

% ============================================================
% ОТВЕТЫ
% ============================================================

\newpage

\thispagestyle{fancy}

\begin{center}
{\LARGE\bfseries\color{darkaccent}
Ответы к тренировочному блоку\par}

\vspace{2mm}

{\color{softgray}
Сверяйтесь после самостоятельной попытки.\par}
\end{center}

\vspace{6mm}

\begin{table}[h]

\centering

\caption{Ответы}

\renewcommand{\arraystretch}{1.45}

\begin{tabularx}{\linewidth}{@{}cX@{}}

\toprule
\textbf{№} & \textbf{Ответ} \\
\midrule

1
&
$AO=8$, \quad $OA_1=4$
\\

2
&
$9$
\\

3
&
$34$
\\

4
&
$32$
\\

5
&
$16$
\\

6
&
$90$
\\

7
&
$CK=7$, \quad $BK=9$
\\

8
&
$S_{AOB}=S_{BOC}=S_{COA}=30$
\\

9
&
$\displaystyle \frac{96\sqrt5}{5}$
\\

10
&
$\displaystyle \frac{216\sqrt5}{5}$
\\

\bottomrule

\end{tabularx}

\end{table}

\vfill

\begin{center}
{\small\color{softgray}
Ключевой маршрут:
теория о медианах
$\to$
площади
$\to$
удвоение медианы
$\to$
параллелограмм
$\to$
подобие и прямоугольные треугольники.
}
\end{center}

\end{document}